\chapter{Objetivos}
\label{cap:objetivos}

\section{Objetivo general}

Determinar si las técnicas de adaptación al dominio mejoran el reconocimiento automático
del habla en español para su uso en accesibilidad educativa, y construir un sistema de
subtitulado en vivo que permita comprobarlo en condiciones de aula.

El objetivo se formula deliberadamente como una pregunta abierta y no como la búsqueda de
una mejora: un resultado negativo bien medido responde igual de bien que uno positivo,
siempre que el procedimiento permita distinguir la ausencia de efecto de la incapacidad
para detectarlo.

\section{Objetivos específicos}

\subsection*{Del núcleo investigador}

\begin{enumerate}
  \item Establecer una línea base reproducible del reconocimiento sin adaptar, sobre varios
        modelos y corpus, con trazabilidad completa de la configuración que produjo cada
        cifra.
  \item Evaluar tres técnicas de adaptación (prompting contextual, post-corrección con un
        modelo de lenguaje local y ajuste fino con LoRA) bajo un diseño común que permita
        compararlas entre sí.
  \item Determinar la potencia estadística necesaria para que la comparación sea capaz de
        detectar los efectos esperables, y verificar que el material empleado la alcanza.
  \item Identificar qué modos de fallo relevantes para accesibilidad no quedan reflejados
        en la métrica habitual, y proponer instrumentos que los expongan.
\end{enumerate}

\subsection*{Del núcleo aplicado}

\begin{enumerate}\setcounter{enumi}{4}
  \item Construir una aplicación de subtitulado en vivo desplegable en la red local de un
        centro, con un emisor y múltiples receptores, en la que todo el cálculo resida en
        el servidor.
  \item Caracterizar experimentalmente el compromiso entre latencia y calidad que impone el
        troceado del audio, y elegir el punto de operación con datos propios.
  \item Cuantificar el margen acústico que requiere el sistema, para poder formular
        requisitos de despliegue verificables.
  \item Mantener una frontera entre investigación y aplicación que permita incorporar la
        técnica ganadora al sistema desplegado sin modificarlo.
\end{enumerate}

\section{Criterios de cumplimiento}

Cada objetivo se considera alcanzado bajo condiciones explícitas, fijadas antes de
ejecutar los experimentos para evitar ajustarlas a los resultados obtenidos:

\begin{table}[h]
\centering
\small
\begin{tabular}{|c|p{10.5cm}|}
\hline
\textbf{Obj.} & \textbf{Se considera cumplido si} \\
\hline
1 & Cualquier cifra publicada puede regenerarse con un solo comando, y el resultado
    registra el commit y el contenido exacto del corpus empleado. \\
\hline
2 & Las tres técnicas se miden sobre el mismo corpus, con la misma decodificación y el
    mismo criterio de significación. \\
\hline
3 & El tamaño de muestra se justifica con un cálculo previo, y las mediciones que no lo
    alcanzan se marcan como no concluyentes en lugar de interpretarse. \\
\hline
4 & Se dispone de al menos una métrica adicional que detecte fallos invisibles al WER, con
    evidencia de que tales fallos existen en el material analizado. \\
\hline
5 & Un cliente conectado a la red del centro visualiza los subtítulos sin ejecutar ningún
    cálculo ni requerir permisos de captura. \\
\hline
6 & La elección del tamaño y la estrategia de segmentación se respalda con un barrido
    propio, no con valores tomados por convención. \\
\hline
7 & Existe una curva de degradación frente a ruido que permita expresar el requisito de
    captación en términos verificables. \\
\hline
8 & La técnica ganadora se incorpora al sistema desplegado sin modificar el código de la
    aplicación. \\
\hline
\end{tabular}
\caption{Criterios de cumplimiento de los objetivos específicos.}
\label{tab:criterios}
\end{table}

\section{Alcance y exclusiones}

Quedan fuera del alcance, y se declaran aquí para evitar que se interpreten como
carencias: la evaluación con usuarios reales del sistema desplegado; la traducción a otros
idiomas o a lengua de signos; el reconocimiento de múltiples hablantes simultáneos; y la
identificación de quién habla en cada momento. El sistema asume un único emisor, que es el
escenario de la clase magistral.
